\section{Hamiltonian Derivation and Noise Analysis}
\label{app:hamiltonian-derivation-noise}
This appendix provides an analytical approximation showing that, to leading order, the dominant effect of flux noise is a fluctuation of the FTT frequency.

The full Hamiltonian of the system shown in Fig.\ \ref{fig:device}(a) is
\begin{widetext}
\begin{align}
\label{Hfull}
H_{\text{full}}(t)=\;&4E_C (n-n_g)^2
-E_{J,\mathrm{eff}}\!\left(\Phi+\delta\Phi(t)\right)
\cos\!\left[{\varphi}-{\varphi_0\!\left(\Phi+\delta\Phi(t)\right)}\right]
+\omega_c c^{\dagger}c
+g' n (c^{\dagger}+c)
+2\Omega'_0\cos(\omega_d t)\,n .
\end{align}
\end{widetext}
Here \(E_C\) is the charging energy, \(n\) is the charge operator, and \(n_g\) is the offset charge. The operator \(c\) annihilates photons in the cavity mode of frequency \(\omega_c\). The charge--charge coupling strength is \(g'\), and the final term describes a charge drive applied to the FTT with amplitude \(\Omega'_0\) and frequency \(\omega_d\).

Both the effective Josephson energy \(E_{J,\mathrm{eff}}\) and the potential phase \(\varphi_0\) depend on flux:
\begin{multline}
    E_{J,\mathrm{eff}}\!\left(\Phi+\delta\Phi(t)\right)
    = (E_{J1}+E_{J2}) \cos\!\left[\frac{\pi\left(\Phi+\delta\Phi(t)\right)}{\Phi_0}\right],\\
    \varphi_0\!\left(\Phi+\delta\Phi(t)\right)
    = d\, \tan\!\left[\frac{\pi\left(\Phi+\delta\Phi(t)\right)}{\Phi_0}\right],
\end{multline}
where \(E_{J1}\) and \(E_{J2}\) are the Josephson energies of the two junctions, \(\Phi_0=h/2e\) is the superconducting flux quantum, and \(d \equiv (E_{J2}-E_{J1})/({E_{J1}+E_{J2}})\) is the SQUID-loop asymmetry.

Flux noise in these systems is typically observed to follow a \(1/f\) characteristic. Accordingly, the power spectral density of these Gaussian fluctuations is
\begin{equation}
    S_{1/f}(\omega) = \int_0^{t} \langle \delta \Phi(0)\, \delta\Phi(\tau) \rangle e^{i\omega \tau}\, d\tau
    = \frac{A^2_0}{|\omega / 2 \pi|},
\end{equation}
where \(A_0\) is the flux-noise spectral-density amplitude.

For small flux fluctuations \(\delta\Phi(t)\ll\Phi_0\), a Taylor expansion of \(H_{\text{full}}\) gives
\begin{equation}
H_{\text{full}}\!\left(\Phi+\delta\Phi(t)\right)
\approx H_{\text{sys}}(\Phi) + \underbrace{\delta\Phi(t)\frac{\partial H_{\text{sys}}}{\partial \Phi}}_{H_{\text{noise}}'} 
+ \mathcal{O}\big((\delta \Phi)^2\big),
\end{equation}
where the system Hamiltonian is evaluated at the static bias point \(\Phi\),
\begin{equation}
\begin{aligned}
\label{system hamiltonian}
H_{\text{sys}}(\Phi) =\; & 4E_C (n-n_g)^2
- E_{J,\mathrm{eff}}(\Phi)
\cos\!\left[{\varphi}-{\varphi_0(\Phi)}\right] \\
& + \omega_c c^{\dagger}c
+ g' n (c^{\dagger}+c)
+ 2\Omega'_0 \cos(\omega_d t)\, n.
\end{aligned}
\end{equation}
The dependence on \(\Phi\) is dropped below for compactness.

In the transmon regime (\(E_C \ll E_{J,\text{eff}}\)), the wavefunction is localized in the potential well, and the Hamiltonian can be expressed using ladder operators:
\begin{equation}
\begin{aligned}
    \varphi - \varphi_0 &= \varphi_{\text{ZPF}} \left( b + b^\dagger \right), \\
    n &= \frac{-i \left( b - b^\dagger \right)}{2 \varphi_{\text{ZPF}}}, \\
    \varphi_{\text{ZPF}} &= \left( \frac{2 E_C}{E_{J,\text{eff}} (\Phi)} \right)^{\frac{1}{4}}.
\end{aligned}
\end{equation}
Expanding the cosine potential to fourth order and applying the rotating wave approximation (RWA) yields
\begin{equation}
\begin{split}
    H_\text{sys}(t) \approx \; & \omega_{b} b^{\dagger} b + \frac{K}{2} b^{\dagger 2} b^{2}+\omega_c c^\dagger c \\
    & + g ( b+ b^{\dagger})( c+ c^{\dagger}) + 2\Omega_0\cos(\omega_{d}t)( b+ b^{\dagger}),
\end{split}
\end{equation}
where the FTT frequency, anharmonicity, renormalized drive amplitude and coupling strength are \(\omega_{b} = \sqrt{8E_{C}E_{J,\text{eff}} } - E_{C}\), \(K = -E_{C}\), \(\Omega_0 = \Omega_0'/2\varphi_{\text{ZPF}}\) and \(g = g'/2\varphi_{\text{ZPF}}\) respectively.

The noise Hamiltonian in the oscillator basis defined at the static flux bias \(\Phi\) is
\begin{equation}
\begin{aligned}
H_{\text{noise}}'(t) &= \delta\Phi(t) \frac{\partial H_{\text{sys}}}{\partial \Phi} \\
&= -\delta\Phi(t)\biggl[ 
\frac{\partial E_{J,\mathrm{eff}}}{\partial\Phi} \cos(\varphi - \varphi_0) \\
&\qquad\qquad + E_{J,\mathrm{eff}} \frac{\partial \varphi_0}{\partial\Phi} \sin(\varphi - \varphi_0) 
\biggr] \\
&\approx -\delta\Phi(t)\biggl[
\frac{\partial E_{J,\mathrm{eff}}}{\partial\Phi}
\left(1-\frac{(\varphi-\varphi_0)^2}{2}\right) \\
&\qquad\qquad + E_{J,\mathrm{eff}}\frac{\partial \varphi_0}{\partial\Phi}(\varphi-\varphi_0)
+ \mathcal{O}\big((\varphi-\varphi_0)^3\big)
\biggr] \\
&\approx -\delta\Phi(t)\biggl[
\frac{\partial E_{J,\mathrm{eff}}}{\partial\Phi}
\left(1-\frac{\varphi_{\mathrm{ZPF}}^2}{2}(b+b^\dagger)^2\right) \\
&\qquad\qquad + E_{J,\mathrm{eff}}\frac{\partial \varphi_0}{\partial\Phi}\,
\varphi_{\mathrm{ZPF}}(b+b^\dagger)
\biggr].
\end{aligned}
\end{equation}
This expression contains a term linear in \((b+b^\dagger)\) and a quadratic contribution proportional to \((b+b^\dagger)^2\). The linear term drives transitions near the oscillator frequency \(\omega_b\), with rates governed by \(S_{1/f}(\omega_b)\). Expanding \((b+b^\dagger)^2 = 2b^\dagger b + 1 + b^2 + b^{\dagger 2}\) shows that the squeezing terms (\(b^2\) and \(b^{\dagger 2}\)) drive transitions near \(2\omega_b\), with rates proportional to \(S_{1/f}(2\omega_b)\). Because \(S_{1/f}(\omega)\) is strongly suppressed at these high frequencies (\(\omega_b/2\pi\sim\) GHz) compared to the low-frequency noise (\(\omega/2\pi \rightarrow 0\) Hz) responsible for pure dephasing, both the linear and squeezing transition channels are neglected. Retaining only the number-conserving contribution (\(\propto b^\dagger b\)) from the quadratic term gives
\begin{equation}
H_{\text{noise}}'(t)\approx
\frac{\partial E_{J,\mathrm{eff}}}{\partial\Phi}\,\varphi_{\mathrm{ZPF}}^2\,
\delta\Phi(t)\; b^\dagger b ,
\end{equation}
up to an additive constant. Using
\begin{equation}
\frac{\partial\omega_b}{\partial\Phi}
=
\frac{\partial\omega_b}{\partial E_{J,\mathrm{eff}}}\frac{\partial E_{J,\mathrm{eff}}}{\partial\Phi}
=
\varphi_{\mathrm{ZPF}}^2\,\frac{\partial E_{J,\mathrm{eff}}}{\partial\Phi},
\label{eq:sensitivity_deriv}
\end{equation}
one obtains
\begin{equation}
H_{\text{noise}}'(t)\approx
\frac{\partial\omega_b}{\partial\Phi}\,\delta\Phi(t)\; b^\dagger b,
\end{equation}
which is the noise Hamiltonian in Eq.~\eqref{H}.